\documentclass{article}
\usepackage{amsmath} 
\usepackage{amsfonts}
\usepackage{amssymb}
\usepackage[parfill]{parskip}
\usepackage{esvect}
\usepackage[thinc]{esdiff}
\usepackage{mathtools}
\usepackage{physics}
\usepackage{amsthm}
\usepackage{xfrac}
\usepackage{enumitem}

\DeclareMathOperator{\GL}{GL}

\begin{document}

\title{MATH 430: HW 4}
\author{Jacob Lockard}
\date{19 February 2026}

\maketitle

\section*{Exercise 2.44}

Let $f : G_1 \to G_2$ be a group isomorphism from the group
$\langle G_1, *_1 \rangle$ to the group $\langle G_2, *_2 \rangle$.
We'll show that $f^{-1} : G_2 \to G_1$ is also a group isomorphism.

Let $a, b \in G_1$ with $f^{-1}(a) = f^{-1}(b)$.
Then, $f\bigl( f^{-1}(a) \bigr) = f\bigl( f^{-1}(b) \bigr)$ and
$a = b$,
so $f^{-1}$ is injective. Let $y \in G_2$. Then $f^{-1}\bigl( f(y) \bigr) = y$, so
$f^{-1}$ is surjective.

Let $x, y \in G_2$. Since $f$ is surjective, there exist $a, b \in {G_1}$
such that $f(a) = x$ and $f(b) = y$. We note that $f^{-1}(x) = a$ and
$f^{-1}(y) = b$.
Since $f$ is a group isomorphism,
we know
\begin{align*}
    f(a *_1 b) &= f(a) *_2 f(b) \\
    f^{-1}\bigl( f(a *_1 b) \bigr) &= f^{-1}\bigl( f(a) *_2 f(b) \bigr) \\
    a *_1 b &= f^{-1}\bigl( f(a) *_2 f(b) \bigr) \\
    f^{-1}(x) *_1 f^{-1}(y) &= f^{-1}\bigl( x *_2 y \bigr) \,.
\end{align*}
So $f^{-1}$ is a group homomorphism. Since $f^{-1}$ is a bijective homomorphism,
we conclude that it is an isomorphism.

\qedsymbol

\section*{Exercise 3.39}

We will show that there does not exist an isomorphism of $U_6$ with
$\mathbb{Z}_6$ in which $\zeta = e^{i(\pi/3)}$ corresponds to $4$.
Let $f : U_6 \to \mathbb{Z}_6$ be a group isomorphism with
$f(\zeta) = 4$. Then we have:
\begin{align*}
    f(\zeta^3) &= 3 f(\zeta) = 3(4) = 2 +_6 4 = 0 \,; \\
    f(\zeta^6) &= 6 f(\zeta) = 6(4) = 3(4) +_6 3(4) = 0 +_6 0 = 0 \,.
\end{align*}
However,
\begin{align*}
    \zeta^3 = \bigl(e^{i(\pi/3)}\bigr)^3 = e^{i \pi} = -1
    \neq 1 = e^{i(2\pi)} = \bigl(e^{i(\pi/3)}\bigr)^6 = \zeta^6 \,.
\end{align*}
So $f$ is not injective, and we conclude that $f$ does not exist.

\qedsymbol

\section*{Exercise 4.12c}

\begin{quote}
    \itshape
    Compute the permutation products.

    \begin{enumerate}[label=\alph*., start=3]
        \item $\bigl[(1,6,7,2)^2(4,5,2,6)^{-1}(1,7,3)\bigr]^{-1}$
    \end{enumerate}
\end{quote}
We have:
\begin{align*}
    (1,6,7,2)^2 = (1, 7)(6, 2) \,; \\
    (4, 5, 2, 6)^{-1} = (6,2,5,4) \,; \\
    (1,7)(6,2)(6,2,5,4)(1,7,3) = (7,3)(2, 5, 4) \,; \\
    [(7,3)(2,5,4)]^{-1} = (3, 7)(4, 5, 2) \,.
\end{align*}

\section*{Exercise 5.23}

\begin{quote}
    \itshape
    Describe all the elements in the cyclic subgroup of
    $\GL(2,\mathbb{R})$ generated by the given $2\times2$ matrix.

    \begin{enumerate}[start=24]
        \item \begin{align*}
            \begin{bmatrix}
                1 & 1 \\
                0 & 1
            \end{bmatrix}
        \end{align*}
    \end{enumerate}
\end{quote}

If $n = 0$, then
\begin{align*}
    \begin{bmatrix}
        1&1\\
        0&1
    \end{bmatrix}^n = \begin{bmatrix}
        1&n\\
        0&1
    \end{bmatrix} \,.
\end{align*}

If for some $n \in \mathbb{N}_0$ it holds that
\begin{align}
    \begin{bmatrix}
        1&1\\
        0&1
    \end{bmatrix}^n = \begin{bmatrix}
        1&n\\
        0&1
    \end{bmatrix} \,,
\end{align} then
\begin{align*}
    \begin{bmatrix}
        1&1\\
        0&1
    \end{bmatrix}^{n+1} = \begin{bmatrix}
        1&1\\
        0&1
    \end{bmatrix}^n \begin{bmatrix}
        1&1\\
        0&1
    \end{bmatrix} = \begin{bmatrix}
        1&n\\0&1
    \end{bmatrix} \begin{bmatrix}
        1&1\\0&1
    \end{bmatrix} = \begin{bmatrix}
        1&n+1\\
        0&1
    \end{bmatrix}\,.
\end{align*}
So by induction, (1) holds for all $n \in \mathbb{N}_0$.

Now let $n \in \mathbb{N}$. Then,
\begin{align*}
    \begin{bmatrix}
        1&n\\0&1
    \end{bmatrix}
    \begin{bmatrix}
        1&-n\\
        0&1
    \end{bmatrix} = I_2 \,,
\end{align*}
so
\begin{align*}
    \begin{bmatrix}
        1&n\\
        0&1
    \end{bmatrix}^{-1} = \begin{bmatrix}
        1&-n\\
        0&1
    \end{bmatrix} \,,
\end{align*}
and
\begin{align*}
    \begin{bmatrix}
        1&1\\
        0&1
    \end{bmatrix}^{-n} = \begin{bmatrix}
        1&-n\\
        0&1
    \end{bmatrix} \,.
\end{align*}
We conclude that the entire cyclic subgroup is all matrices of the form
\begin{align*}
    \begin{bmatrix}
        1&n\\
        0&1
    \end{bmatrix} \,,
\end{align*}
where $n \in \mathbb{Z}$.

\end{document}